\documentclass[11pt]{article}
\usepackage[a4paper,margin=1in]{geometry}
\usepackage{amsmath,amssymb,amsthm,mathtools}
\usepackage{graphicx}
\usepackage[hidelinks]{hyperref}
\usepackage{microtype}
\usepackage[T1]{fontenc}
\usepackage{lmodern}
\setlength{\emergencystretch}{3em}

\newtheorem{theorem}{Theorem}
\newtheorem{lemma}{Lemma}
\newtheorem{remark}{Remark}
\newcommand{\C}{\mathbb C}
\newcommand{\cE}{\mathcal E}
\newcommand{\cO}{\mathcal O}
\newcommand{\cC}{\mathfrak c}
\newcommand{\epsprod}{\varepsilon}

\title{A Product Counterexample to Problem 5.14(iii)\\
\large Strict admissibility does not characterize the three classes in Theorem 4.3}
\author{Candidate solution packet for arXiv:1912.03659}
\date{27 August 2026}

\begin{document}
\maketitle

\noindent\textbf{Status.}
Candidate full negative answer to Problem 5.14(iii), likely valid and awaiting
human review.  The result answers only part (iii) of the four-part problem.

\section{Source question}

The source is Karsten Kruse, \emph{Vector-valued Fourier hyperfunctions and
boundary values}, arXiv:1912.03659v2, 5 November 2025
\cite{Kruse}.  On source PDF page 46, Problem 5.14(iii) asks:

\begin{quote}
``Is strict admissibility of a sequentially complete $E$ equivalent to
belonging to the classes of spaces from Theorem 4.3?''
\end{quote}

Theorem 4.3 gives the following three sufficient classes: (a) complex
Fr\'echet spaces; (b) strong duals $F'_b$ of complex Fr\'echet spaces $F$
satisfying $(DN)$; and (c) complex ultrabornological PLS-spaces satisfying
$(PA)$ \cite[Theorem 4.3]{Kruse}.

\begin{figure}[ht]
  \centering
  \includegraphics[width=\textwidth]{figures/open_problem_crop.png}
  \caption{Problem 5.14 in the source paper, source PDF page 46.}
\end{figure}

\section{Definitions used}

A complex locally convex Hausdorff space $E$ is \emph{admissible} if, for
every compact $K\subset\overline{\mathbb R}$, the Cauchy--Riemann operator
\[
 \overline\partial:\cE^{\exp}(\overline\C\setminus K,E)
 \longrightarrow \cE^{\exp}(\overline\C\setminus K,E)
\]
is surjective.  It is \emph{strictly admissible} if it is admissible and
\[
 \overline\partial:C^\infty(U,E)\longrightarrow C^\infty(U,E)
\]
is surjective for every open $U\subset\C$ \cite[Definition 4.1]{Kruse}.

Write $\cC=|\mathbb R|=|\C|$ for the cardinality of the continuum.  Let $I$
be any set with $|I|=2^{\cC}$ and equip
\[
                         E:=\C^I
\]
with its product locally convex topology.

\section{Proof intuition}

The analytic property in the question is completely coordinatewise for a
product.  A smooth map into $\C^I$ is precisely a family of scalar smooth
maps, and every continuous seminorm on the product sees only finitely many
coordinates.  Scalar $\overline\partial$-solutions can therefore be chosen
independently in every coordinate without losing any weighted estimate.  This
makes $\C^I$ strictly admissible for an arbitrary index set $I$.

The rest of the argument chooses $I$ deliberately large and separates the
resulting product from the three classes by elementary invariants.  It is not
metrizable; it has no countable fundamental family of bounded sets, while a
strong dual of a Fr\'echet space does; and its underlying cardinality is larger
than the continuum, while every PLS-space has cardinality at most the
continuum.  Completeness is preserved by arbitrary products.

\section{Formal result}

\begin{lemma}[Products preserve strict admissibility]\label{lem:product}
Let $(E_j)_{j\in J}$ be any family of strictly admissible complex locally
convex Hausdorff spaces.  Then $\prod_{j\in J}E_j$, with the product topology,
is strictly admissible.  If every $E_j$ is complete, then the product is
complete.
\end{lemma}

\begin{proof}
Put $X=\prod_{j\in J}E_j$ and denote the coordinate projections by $\pi_j$.
A map into a locally convex product is smooth exactly when all its coordinate
maps are smooth, with derivatives computed coordinatewise.  Moreover, every
continuous seminorm $p$ on $X$ is dominated by
\[
             C\max_{j\in F}p_j\circ\pi_j
\]
for some finite $F\subset J$, some continuous seminorms $p_j$ on $E_j$, and
some $C>0$.

Fix a compact $K\subset\overline{\mathbb R}$ and
$f\in\cE^{\exp}(\overline\C\setminus K,X)$.  For every $j$, the coordinate
$f_j=\pi_j\circ f$ belongs to
$\cE^{\exp}(\overline\C\setminus K,E_j)$.  Admissibility of $E_j$ gives
$u_j$ in that same space with $\overline\partial u_j=f_j$.  Define
$u=(u_j)_{j\in J}$.  Coordinatewise smoothness makes $u$ smooth.  The finite
coordinate domination of every continuous seminorm shows, directly in each
defining weighted seminorm, that
$u\in\cE^{\exp}(\overline\C\setminus K,X)$.  Hence
$\overline\partial u=f$.

The identical coordinatewise construction, without weights, solves
$\overline\partial u=f$ for every $f\in C^\infty(U,X)$ and every open
$U\subset\C$.  Thus $X$ is strictly admissible.  Finally, an arbitrary
product of complete Hausdorff uniform spaces is complete, so the last claim
follows.
\end{proof}

\begin{lemma}[Two size obstructions]\label{lem:obstructions}
Let $I$ be uncountable.
\begin{enumerate}
\item The product $\C^I$ is not metrizable and has no countable fundamental
family of bounded sets.
\item Every PLS-space has underlying cardinality at most $\cC$.
\end{enumerate}
\end{lemma}

\begin{proof}
For the first assertion, suppose $(V_n)$ were a countable neighborhood base at
zero.  Each product neighborhood restricts only finitely many coordinates.
The union of those finite coordinate sets is countable, so one may choose
$i\in I$ outside it.  Every $V_n$ then contains the whole $i$-th coordinate
line, whereas the neighborhood $\{x:|x_i|<1\}$ does not.  This is a
contradiction.

A subset of $\C^I$ is bounded exactly when it is bounded in every coordinate.
Given any sequence $(B_n)$ of bounded sets, choose distinct $i_n\in I$, put
\[
       M_n:=\sup_{x\in B_n}|x_{i_n}|<\infty,
\]
and define $x\in\C^I$ by $x_{i_n}=M_n+1$ and $x_i=0$ on all other
coordinates.  The singleton $\{x\}$ is bounded, but $x\notin B_n$ for every
$n$.  Hence no countable sequence of bounded sets is cofinal.

For the second assertion, recall that a DFS-space is the strong dual of a
Fr\'echet--Schwartz space.  Every Fr\'echet--Schwartz space is separable: in a
Banach-step representation the relevant linking maps are compact with dense
range, so each Banach step is separable, and the Fr\'echet space embeds into a
countable product of separable metrizable spaces.  A continuous functional on
a separable Hausdorff topological vector space is determined by its values on
a countable dense set.  Therefore every DFS-space has cardinality at most
$\cC$.  A PLS-space is a projective limit of a sequence of DFS-spaces, hence
is a subset of their countable product and has cardinality at most
$\cC^{\aleph_0}=\cC$.
\end{proof}

\begin{theorem}[Negative answer to Problem 5.14(iii)]\label{thm:main}
Strict admissibility of a sequentially complete complex locally convex
Hausdorff space is not equivalent to membership in the three classes of
Theorem 4.3 of \cite{Kruse}.

More explicitly, if $|I|=2^{\cC}$, then $E=\C^I$ is complete and strictly
admissible, but it is neither
\begin{enumerate}
\item a Fr\'echet space,
\item a strong dual $F'_b$ of a Fr\'echet space $F$ satisfying $(DN)$, nor
\item an ultrabornological PLS-space satisfying $(PA)$.
\end{enumerate}
\end{theorem}

\begin{proof}
The scalar field $\C$ is a complete strictly admissible space; this is also the
one-dimensional instance of Theorem 4.3(a) in \cite{Kruse}.  Lemma
\ref{lem:product} therefore shows that $E=\C^I$ is complete, hence
sequentially complete, and strictly admissible.

Because $I$ is uncountable, Lemma \ref{lem:obstructions}(1) shows that $E$ is
not metrizable, so it is not Fr\'echet.

Suppose next that $E$ were topologically isomorphic to $F'_b$ for some
Fr\'echet space $F$ (the condition $(DN)$ is not needed here).  Let $(U_n)$ be
a countable zero-neighborhood base in $F$.  Every bounded subset $B$ of the
strong dual is pointwise bounded and hence equicontinuous by Banach--Steinhaus,
because $F$ is barrelled.  Thus $B\subset U^\circ$ for some zero-neighborhood
$U$; choosing $n$ with $U_n\subset U$ gives
$B\subset U^\circ\subset U_n^\circ$.  Consequently $(U_n^\circ)$ is a
countable fundamental family of bounded sets in $F'_b$.  This contradicts
Lemma \ref{lem:obstructions}(1).

Finally,
\[
 |E|=|\C|^{|I|}=\cC^{,2^{\cC}}
      =(2^{\aleph_0})^{2^{\cC}}=2^{2^{\cC}}>\cC.
\]
Lemma \ref{lem:obstructions}(2) therefore rules out even that $E$ is a
PLS-space, a fortiori that it belongs to class (c).  This excludes all three
classes and proves the result.
\end{proof}

\section{Verification notes}

\begin{itemize}
\item \textbf{Hypotheses.}  The example is a complex locally convex Hausdorff
space and is complete, not merely sequentially complete.
\item \textbf{Analytic step.}  Both weighted and unweighted
$\overline\partial$ equations are checked coordinatewise.  The only topology
fact used is that continuous seminorms on a product depend on finitely many
coordinates.
\item \textbf{Class (b).}  The obstruction is stronger than required: the
example is not the strong dual of any Fr\'echet space, regardless of $(DN)$.
\item \textbf{Class (c).}  The cardinal argument rules out the PLS property
itself, before $(PA)$ or ultrabornologicity is considered.
\item \textbf{Set theory.}  The choice $|I|=2^{\cC}$ and the strict inequality
$2^{2^{\cC}}>\cC$ use only Cantor's theorem and are valid in ZFC.
\end{itemize}

\section{Scope and novelty check}

This settles the literal form of Problem 5.14(iii) negatively.  It does not
address Problems 5.14(i), (ii), or (iv).  The counterexample is intentionally
very large and nonmetrizable.  If the intended question was restricted to an
unstated smaller category of ``natural'' locally convex spaces, that narrower
variant remains open and should be distinguished from the printed statement.

A bounded novelty check was performed on 27 August 2026.  The four local run
indexes were searched for arXiv:1912.03659 and the phrases ``strict
admissibility'', ``admissible'', ``Fourier hyperfunctions'', and ``product''.
Web/arXiv searches used the exact wording of Problem 5.14(iii), the paper title
and author, and the phrases ``strictly admissible product'' and ``classes of
spaces from Theorem 4.3''.  The current arXiv v2 source, dated 5 November 2025,
still states the problem.  No separate paper or packet asserting this product
counterexample or answering Problem 5.14(iii) was found.  This is a bounded,
not exhaustive, novelty search; novelty confidence is therefore moderate.

\section{Human-review recommendation}

Promote as a full negative answer to Problem 5.14(iii) if a reviewer confirms
(1) that the problem is intended literally over all sequentially complete
complex locally convex Hausdorff spaces, and (2) the standard cardinal bound
$|X|\leq\cC$ for PLS-spaces.  The coordinatewise strict-admissibility argument
and the strong-dual bounded-set obstruction are otherwise elementary.

\begin{thebibliography}{9}
\bibitem{Kruse}
K.~Kruse,
\emph{Vector-valued Fourier hyperfunctions and boundary values},
arXiv:1912.03659v2 (5 November 2025),
\url{https://arxiv.org/abs/1912.03659};
published in \emph{Kyoto Journal of Mathematics},
\href{https://doi.org/10.1215/21562261-2024-0030}{doi:10.1215/21562261-2024-0030}.
\end{thebibliography}

\end{document}
